%% THERMODYNAMIQUE %%
%
\section{Composantes du vent \`a partir des vitesses "avion" dans le rep\`ere avion}

\itemstyle{ Op\'erateur~: } {\sf CalcVent}
\index{CalcVent}

%%%%%OBJET%%%%%
\itemstyle{ Objet~: }
Calcul des 3 composantes du vecteur vent � partir de mesures faites
dans un rep�re li� � une plate-forme mobile (avion, bateau, ...).

Cet op�rateur est adapt� au cas o� la mesure du vent relatif est exprim�e
en module, angles d'incidence et d�rapage, et n'est donc utilisable
que lorsque le vent vient uniquement de l'avant de la plate-forme mobile
(l'incidence n'est pas d�finie lorsque le d�rapage est de +/- $\pi/2$).

Le calcul est effectu\'e dans le rep\`ere li\'e au mobile, puis les composantes 
sont transform\'ees vers le rep\`ere g\'eographique terrestre local
(ENU : axes orient\'es vers l'Est, le Nord, et verticalement).

Cet op\'erateur ne traite que le cas particulier o� les points
de mesure de la vitesse du vent relatif (l'air par rapport au mobile)
et le point de mesure de la vitesse par rapport au sol sont
align�s sur l'axe longitudinal.

%%%%%ARGUMENTS%%%%
\bigskip
%%%%%ENTREES%%%%
\begin{tabarg}
\multicolumn{3}{l}{\titretab{ Entr\'ees}}\\
\hline
$L$ & Coeff., float& Distance alg\'ebrique entre le point de mesure de la vitesse 
de l'air par rapport \`a l'avion, et le point de mesure de la 
vitesse avion par rapport au sol ($m$)\\
$V_p$ & Vecteur, {\tt float} & Vitesse propre ($m/s$) \\
$\alpha$ & Vecteur, {\tt float} & Angle d'attaque, positif lorsque le vent relatif vient de sous l'avion ($radians$) \\
$\beta$ & Vecteur, {\tt float} & Angle de d\'erapage, positif lorsque le 
vent relatif vient de tribord ($radians$)\\
$V_x$ & Vecteur, {\tt float} & Vitesse de l'avion par rapport au sol suivant 
l'axe longitudinal de l'avion dirig\'e vers l'avant ($m/s$)\\
$V_y$ & Vecteur, {\tt float} & Vitesse de l'avion par rapport au sol suivant l'axe 
transversal de l'avion dirig\'e de babord vers tribord ($m/s$)\\ 
$V_z$ & Vecteur, {\tt float} & Vitesse de l'avion par rapport au sol suivant l'axe 
vertical de l'avion dirig\'e vers le bas ($m/s$)\\
$Vr_y$ & Vecteur, {\tt float} & Vitesse de rotation suivant l'axe Y, positif quand $\theta$ cro\^{\i}t($rad/s$) \\
$Vr_z$ & Vecteur, {\tt float} & Vitesse de rotation suivant l'axe Z, positif quand $\psi$ cro\^{\i}t ($rad/s$) \\
$\theta$ & Vecteur, {\tt float} & Angle de tangage de l'avion, positif quand
l'avion est cabr\'e ($radians$)\\
$\phi$ & Vecteur, {\tt float} & Angle de roulis de l'avion, positif quand
l'aile tribord est enfonc\'ee ($radians$)\\
$\psi$ & Vecteur, {\tt float} & Cap g\'eographique de l'avion ($radians$) \\
\end{tabarg}

\medskip
%%%%%SORTIES%%%%
\begin{tabarg}
\multicolumn{3}{l}{\titretab{ Sorties}}\\
\hline
$U$ & Vecteur, {\tt float} & Composante zonale du vent (Ouest en Est) ($m/s$)\\
$V$ & Vecteur, {\tt float} & Composante m\'eridienne du vent (Sud au Nord) ($m/s$)\\
$W$ & Vecteur, {\tt float} & Composante verticale du vent (Bas en Haut) ($m/s$)\\
\end{tabarg}



%%%%%FORMULE%%%%
\itemstyle{Formule}

Dans un premier temps, le vent est calcul\'e dans le 
rep\`ere avion. Ses 3 composantes sont not\'ees $U_a$, $V_a$,
et $W_a$. 

\begin{align*}
U_a & = & V_x - V_p.D^{-1} \\
V_a & = & V_y - V_p.D^{-1}.\tan\beta + L.Vr_z\\
W_a & = & V_z - V_p.D^{-1}.\tan\alpha - L.Vr_y
\end{align*}

avec~: $ D = \sqrt{1+\tan^2\alpha+\tan^2\beta}$\\

Puis, un changement de rep\`ere est effectu\'e pour obtenir le 
vent dans le rep\`ere g\'eographique.\\
\\
$U = U_a.\sin\psi.\cos\theta +$
        $V_a.(\cos\psi.\cos\phi+\sin\psi.\sin\theta.\sin\phi) +$
        $W_a.(\sin\psi.\sin\theta.\cos\phi - \cos\psi.\sin\phi)$ \\   
$V = U_a.\cos\psi.\cos\theta -$
        $V_a.(\sin\psi.\cos\phi - \cos\psi.\sin\theta.\sin\phi) +$
        $W_a.(\cos\psi.\sin\theta.\cos\phi + \sin\psi.\sin\phi)$\\
$W = U_a.\sin\theta -$ 
        $V_a.\cos\theta.\sin\phi -$
        $W_a.\cos\theta.\cos\phi$
\\



%%%%%COMMENTAIRE%%%%%
\itemstyle{Commentaire~:}
Pour le cas g�n�ral, voir l'op�rateur CalcVentNA.

%%%%%AUTEUR%%%%%
\itemstyle{Auteur~:}
CNRM/GMEI/TRAMM

%%%%%REFERENCES%%%%%
\itemstyle{R\'ef\'erences~:}
Bulletin NCAR-RAF Nr 23 \cite{RAF-23}.
